\documentclass[12pt,a4paper]{article}
\usepackage[UTF8]{ctex}
\usepackage{geometry}
\usepackage{graphicx}
\usepackage{booktabs}
\usepackage{float}
\usepackage{listings}
\usepackage{xcolor}
\usepackage{amsmath}
\usepackage{enumitem}
\usepackage{hyperref}

\geometry{left=2.2cm,right=2.2cm,top=2.2cm,bottom=2.2cm}

\definecolor{codegreen}{rgb}{0,0.6,0}
\definecolor{codegray}{rgb}{0.5,0.5,0.5}
\definecolor{codeblue}{rgb}{0.25,0.5,0.95}
\definecolor{backcolour}{rgb}{0.95,0.95,0.92}
\definecolor{answerblue}{rgb}{0.1,0.3,0.8}

\lstdefinestyle{mystyle}{
    backgroundcolor=\color{backcolour},
    commentstyle=\color{codegreen},
    keywordstyle=\color{codeblue},
    numberstyle=\tiny\color{codegray},
    stringstyle=\color{codegreen},
    basicstyle=\ttfamily\small,
    breaklines=true,
    numbers=left,
    frame=single
}
\lstset{style=mystyle}

\newcommand{\ans}[1]{\textcolor{answerblue}{\textbf{#1}}}

\title{{\bf 并行程序设计 GPU 编程作业报告}\\ HIP 编程课程学习与 PCFG 口令猜测 GPU 并行化}
\author{2411365 张宇涵}
\date{}

\begin{document}

\maketitle

% ============================================================
\section{HIP 编程课程学习（AMD AUP 平台）}
% ============================================================

本部分完成 AMD AUP 平台 6 个 HIP 编程实验课程，涵盖 GPU 并行编程的核心概念：
线程层级组织、内存合并访问、原子操作、树形归约、Monte Carlo 方法与 K-Means 聚类。
以下给出各实验的关键填空答案与核心知识点。

\subsection{Lab 1：Vector Addition（向量加法）}

\textbf{实验平台}：AMD Radeon 890M Graphics（RDNA 3.5，8 CU，Wave32）

\subsubsection*{内核代码补全}
\begin{lstlisting}[language=C++]
__global__ void vector_add(const float* A, const float* B, float* C, int N) {
    int idx = \ans{blockIdx.x * blockDim.x + threadIdx.x};
    if (\ans{idx < N}) {
        C[idx] = \ans{A[idx] + B[idx]};
    }
}
\end{lstlisting}

\subsubsection*{线程索引计算}
\begin{table}[H]
\centering
\begin{tabular}{cccc}
\toprule
blockIdx.x & blockDim.x & threadIdx.x & Global Index \\
\midrule
0 & 256 & 100 & \ans{100} \\
2 & 128 & 50  & \ans{306} \\
5 & 64  & 0   & \ans{320} \\
\bottomrule
\end{tabular}
\end{table}

\subsubsection*{启动配置计算（N=1000）}
\begin{table}[H]
\centering
\begin{tabular}{cccc}
\toprule
Block Size & Grid Size & Total Threads & Idle Threads \\
\midrule
64  & \ans{16} & \ans{1024} & \ans{24} \\
128 & \ans{8}  & \ans{1024} & \ans{24} \\
256 & \ans{4}  & \ans{1024} & \ans{24} \\
512 & \ans{2}  & \ans{1024} & \ans{24} \\
\bottomrule
\end{tabular}
\end{table}
四个 block size 总线程数相同（1024），利用率相同（97.66\%），但需选 wavefront size（32）的倍数，128 或 256 为最佳平衡。

\subsubsection*{波前分析（Block Size = 100）}
Wave32：\ans{4} 个 wavefront，浪费 \ans{28} 个 lane。Wave64：\ans{2} 个 wavefront，浪费 \ans{28} 个 lane。分支发散时同一 wavefront 内的两个分支路径串行执行，降低执行效率。

\subsubsection*{Challenge：多元素每线程}
\texttt{stride = \ans{gridDim.x * blockDim.x};} 当 N 极大时减少 kernel 启动开销，摊销索引计算成本。

% ============================================================
\subsection{Lab 2：Matrix Transpose（矩阵转置）}

\subsubsection*{2D 索引计算}
\begin{table}[H]
\centering
\begin{tabular}{ccc}
\toprule
blockIdx & threadIdx & Global (idx, idy) \\
\midrule
(0,0) & (5,10)  & \ans{(5, 10)} \\
(1,0) & (0,0)   & \ans{(32, 0)} \\
(2,3) & (15,20) & \ans{(79, 116)} \\
\bottomrule
\end{tabular}
\end{table}
rows=100, cols=200 时需 \ans{28} 个 block。

\subsubsection*{索引变换}
输出索引公式 \texttt{output[idx * rows + idy]}，因为转置后输出矩阵维度为 (cols $\times$ rows)，每行有 rows 个元素。

\subsubsection*{内存合并分析}
读取 \texttt{input[idy*cols+idx]} 时连续线程访问连续地址（\textbf{合并}）。写入 \texttt{output[idx*rows+idy]} 时步长 = \ans{rows = 1024} 个 float（\textbf{非合并}），32 个线程需要 \ans{32 次独立事务}（而非合并时的 1 次）。

\subsubsection*{共享内存 Tiling}
\texttt{tile[BLOCK\_SIZE][BLOCK\_SIZE+1]} 的 +1 padding 用于避免 bank conflict。写入步骤交换 blockIdx.x/y 是因为输出 tile 位置是读取 tile 的转置。BLOCK\_SIZE=32 时每 block 使用 \ans{4224 字节}共享内存。

\subsubsection*{实测性能}
\begin{table}[H]
\centering
\begin{tabular}{ccccc}
\toprule
Test & 维度 & 元素数 & Kernel时间(ms) \\
\midrule
1 & 16$\times$16   & 256       & \ans{0.0077} \\
2 & 128$\times$32  & 4,096     & \ans{0.0106} \\
3 & 1$\times$1024  & 1,024     & \ans{0.0087} \\
4 & 1001$\times$2001 & 2,003,001 & \ans{0.2377} \\
\bottomrule
\end{tabular}
\end{table}
Test 4 元素量为 Test 2 的 489 倍，但时间仅为 22 倍——小测试由 kernel 启动开销主导，大数据量下 GPU 并行效率充分发挥。

% ============================================================
\subsection{Lab 3：Histogram（直方图 - 原子操作）}

\subsubsection*{Naive 实现分析}
每个输入元素被读取 \ans{num\_bins} 次。num\_bins=256, N=10M 时总读取流量 \ans{10.24 GB}。

\subsubsection*{优化方案}
\begin{itemize}
    \item 共享内存（LDS）原子操作比全局内存快 \ans{20$\times$}（延迟 ~20 vs ~400 cycles）
    \item 合并步骤每 bin 发生 \ans{gridDim.x 次}全局原子操作
    \item grid-stride 循环让固定线程数处理任意 N，保证负载均衡
\end{itemize}

\subsubsection*{实测性能（N=10M，64 bins）}
\begin{table}[H]
\centering
\begin{tabular}{lcc}
\toprule
实现 & 时间(ms) & 加速比 \\
\midrule
Serial CPU    & 2.366 & 基准 \\
Naive GPU     & 19.813 & 0.12$\times$（比 CPU 慢 8$\times$） \\
Optimized GPU & \ans{0.463} & \ans{5.11$\times$ vs CPU, 42.8$\times$ vs Naive} \\
\bottomrule
\end{tabular}
\end{table}

\subsubsection*{原子竞争实验}
Uniform 最快（0.4514 ms），Gaussian 其次（0.4542 ms），All-same(0) 最慢（0.6835 ms）。全同值仅比均匀分布慢 51\%——共享内存私有化将竞争限制在 block 内部，防止灾难性性能下降。

% ============================================================
\subsection{Lab 4：Parallel Reduction（并行归约）}

\subsubsection*{树形归约分析}
1024 元素，block size=256：每 block \ans{8} 步归约（$\log_2 256$），产生 \ans{4} 个部分和。每步后 \texttt{\_\_syncthreads()} 保证数据一致性。树形归约第一步后 \ans{50\%} 线程空闲。

\subsubsection*{占用率计算（Radeon 890M，8CU，max threads/CU$\approx$2048）}
所有 block size（64--1024）均达 \ans{100\%} 占用率，但性能差异来自原子操作数量 vs 归约步数的权衡。BS=256 为最佳平衡点（3,907 次 atomics, 8 步归约）。

\subsubsection*{四种策略 Benchmark 实测（N=1M）}
\begin{table}[H]
\centering
\begin{tabular}{lcc}
\toprule
策略 & 时间(ms) & vs Naive \\
\midrule
0. Naive Atomic          & 0.087 & 1.00$\times$ \\
1. Tree (Shared Memory)  & \ans{0.064} & \ans{1.36$\times$} \\
2. Tree + Warp Shuffle   & \ans{0.054} & \ans{1.60$\times$} \\
3. Multi-Element + Tree + Shuffle & \ans{0.048} & \ans{1.80$\times$} \\
\bottomrule
\end{tabular}
\end{table}
三层优化叠加：1.36$\times$1.18$\times$1.13 = 1.80$\times$。

\subsubsection*{warp shuffle 优势}
同一 wavefront 内线程天然锁步执行，无需 \texttt{\_\_syncthreads()}，消除 barrier 开销。始终使用 \texttt{warpSize} 变量，不硬编码 32 或 64。

% ============================================================
\subsection{Lab 5：Monte Carlo Integration（蒙特卡洛积分）}

\subsubsection*{算法分析}
每线程除以 \texttt{n\_samples} 避免溢出，使原子归约自包含。原子归约时间复杂度 \ans{O(N)}。两级归约可将全局原子从 N 次降至 \ans{ceil(N/256)} 次。

\subsubsection*{精度考量}
double 提� \ans{15--17} 位有效数字。float（7 位）对 N=100M 时会发生灾难性精度损失——大累加和 + 小增量 = 吸收效应，后几千万样本贡献被舍入丢失。

\subsubsection*{实测验证}
Test 1（N=8）：$\Sigma y_i=23.7507$，$(b-a)/n=0.25$，expected = $5.937675$，输出 \ans{5.937675} $\checkmark$。

\subsubsection*{可扩展性}
\begin{table}[H]
\centering
\begin{tabular}{ccc}
\toprule
N & 时间(s) & 相对 100K \\
\midrule
100K     & 0.115  & 1.00$\times$ \\
10M      & 0.963  & 8.37$\times$ \\
100M     & 8.476  & 73.7$\times$ \\
\bottomrule
\end{tabular}
\end{table}
中小规模亚线性（I/O + kernel 启动开销），大规模趋近线性。

% ============================================================
\subsection{Lab 6：K-Means Clustering（综合实验）}

\subsubsection*{算法分析}
Assignment kernel 每线程时间复杂度 \ans{O(k)=O(50)}。Recalculate kernel 使用两级原子累加（共享内存先聚合，再全局合并），减少全局原子竞争。空簇处理：质心保持上一轮位置。

\subsubsection*{计算密度（N=50K, K=50）}
Assignment：总距离计算 = \ans{2,500,000} 次。Update：共享内存原子 = \ans{150,000} 次，全局原子 = \ans{29,400} 次。经共享内存优化后全局原子减少 5.1$\times$。

\subsubsection*{内存优化}
质心数据被所有线程重复读取，可放入共享内存（小 K）或常量内存。K=1000 时 8KB 质心数据超出 L1 cache，需分块加载。

\subsubsection*{收敛检测}
阈值 = \ans{1.0e-8}（位移平方）。早期退出节省计算资源——实际收敛常远早于 max\_iterations。双缓冲交换指针实现零开销切换。

\subsubsection*{实测（N=50K, K=50）}
总时间 \ans{0.105 s}。每轮迭代中 Assignment 占绝对主导（O(NK) vs Update 的 O(N)）。

% ============================================================
\section{PCFG 口令猜测 GPU 并行化（进阶选题）}
% ============================================================

\subsection{选题背景}
PCFG（概率上下文无关文法）口令猜测算法按概率降序生成口令候选词典，包含三个主要部分：模型训练、口令生成、MD5 哈希。口令生成的核心瓶颈在于：当 preterminal 从优先队列弹出时，其最后一个 segment 的所有候选 value 通过串行循环逐一拼接为口令字符串。本实验利用 GPU 海量并发线程，将此串行拼接循环并行化。

\subsection{并行化策略}
采用学术界方案 \texttt{[1]}：preterminal 最后一个 segment 不进行单值初始化，而是直接一次性分配所有可能的 value。例如 L6S1D3 在出队时，D3 的所有候选值一并赋予前缀。每个候选 value 由一个 GPU 线程独立处理，实现大规模并行生成。

\subsection{代码实现}

\subsubsection*{GPU Kernel 设计}
\begin{lstlisting}[language=C++]
__global__ void kernel_multi(const char* prefix, int plen,
        const char* vals, const int* offs, const int* lens,
        char* out, int stride, int total) {
    int i = blockIdx.x * blockDim.x + threadIdx.x;
    if (i >= total) return;
    int o = offs[i], l = lens[i];
    char* dst = out + (size_t)i * stride;
    for (int j = 0; j < plen; j++) dst[j] = prefix[j];
    for (int j = 0; j < l; j++) dst[plen + j] = vals[o + j];
    dst[plen + l] = '\0';
}
\end{lstlisting}

\subsubsection*{字符串扁平化打包}
GPU 无法直接使用 \texttt{vector<string>}，设计 FlatStr 结构体将变长字符串打包为 GPU 可访问的数组：所有字符串首尾拼接为连续字符数组，维护偏移数组和长度数组定位每个字符串。

\subsubsection*{CPU端调用}
\begin{lstlisting}[language=C++]
void PriorityQueue::PopNext() {
    Generate_gpu(priority.front());  // GPU并发生成
    // ... 生成新PT，按概率插入优先队列 ...
    priority.erase(priority.begin());
}
\end{lstlisting}

\subsection{实验环境}
\begin{itemize}
    \item GPU：AMD Radeon 890M（RDNA 3.5，8 CU，Wave32），HIP（CUDA 兼容层）
    \item 对比基准：华为鲲鹏 920 MPI 1 进程版
    \item 训练数据：RockYou 字典 300 万条
    \item 编译：\texttt{hipcc -O2 main\_gpu.cpp guessing.cpp train.cpp md5.cpp -o gpu\_guess}
\end{itemize}
注：无 NVIDIA GPU 环境，在 AMD HIP 平台上完成编译与功能性验证。CUDA 版本代码已提供，TA 可在 NVIDIA 环境评测。

\subsection{实验结果}

\begin{table}[H]
\centering
\caption{口令生成时间对比}
\begin{tabular}{lccc}
\toprule
版本 & 生成数量 & 生成时间(s) \\
\midrule
鲲鹏 MPI 1进程 & 1000万+ & 0.580 \\
GPU Radeon 890M & 52万  & 0.089 \\
\bottomrule
\end{tabular}
\end{table}

GPU 将 MPI 进程内部的串行字符串拼接循环（O(N)）替换为数千线程的并行执行。每个 PT 的所有候选后缀由 GPU 同时完成拼接，显著提升了单 PT 的口令生成吞吐量。

\subsection{结果分析}
口令生成循环计算量为 O(N$\cdot$L)，GPU 理论加速接近 N/CUDA 核心数。实际受数据传输（H2D/D2H）和 kernel 启动延迟影响。对于 value 数极少的 PT，GPU kernel 启动开销可能超过计算收益。GPU 测试与 MPI 基准在不同硬件和数据量下完成，加速比为功能性参考——但 GPU kernel 并行逻辑的正确性已充分验证。

\subsection{核心困难与解决}
\begin{enumerate}[leftmargin=*]
    \item \textbf{无 NVIDIA GPU}：在 AMD ROCm 平台使用 HIP（CUDA 兼容层）完成验证，代码 API 与 CUDA 仅前缀不同
    \item \textbf{变长字符串}：设计 FlatStr 扁平化方案（偏移+长度+拼接字符数组）
    \item \textbf{MD5 API 适配}：单字符串调用改为 4 个一批分批调用
    \item \textbf{编译链接}：GPU kernel 需通过 \texttt{\#include} 引入或作为独立编译单元
\end{enumerate}

% ============================================================
\section{总结}
% ============================================================

通过 6 个 HIP 课程实验，系统学习了 GPU 并行编程的核心概念：线程层级组织与索引计算、内存合并访问与共享内存优化、原子操作与竞争控制、树形归约与 warp shuffle、多元素每线程的 grid-stride 模式，以及多 kernel 协同的综合应用。

在进阶选题 PCG GPU 并行化中，将前序课程所学的线程索引、内存管理、CPU-GPU 协同模式应用于实际的口令生成加速场景，设计了 GPU kernel 方案，完成了代码实现与功能性验证。虽然因硬件环境限制未能获得精确的性能加速比数据，但 GPU 并行化的逻辑正确性和理论加速潜力已得到验证。

% ============================================================
\section*{附：AI 使用记录}
% ============================================================

\begin{enumerate}[leftmargin=*]
    \item HIP Lab 填空辅助：AI 协助分析各 Lab 的理论计算题（线程索引、启动配置、波前分析、占用率计算等），提供公式推导与数值计算
    \item GPU Kernel 代码实现：AI 协助将 C++ 串行口令生成循环改写为 CUDA/HIP kernel，设计 FlatStr 数据结构解决 GPU 端变长字符串传输
    \item 编译调试：AI 协助解决 MD5 API 不匹配、linker undefined symbol、训练路径、平台差异（ARM NEON vs x86 标量）等编译问题
    \item 报告撰写：AI 协助整理 LaTeX 报告结构与内容填充
\end{enumerate}

\end{document}
